% Part: turing-machines
% Chapter: machines-computations
% Section: variants

\documentclass[../../../include/open-logic-section]{subfiles}

\begin{document}

% SEGMENT OLP-0262-S01

\olfileid{tur}{mac}{var}
\olsection{டியூரிங் பொறிகளின் மாறுவகைகள்}

டியூரிங் பொறிகளை வரையறுக்க உண்மையில் பல வழிகள் உள்ளன; அவற்றுள் நமது
வரையறை ஒரு வழி மட்டுமே. சில வகைகளில் நமது வரையறை மற்றவற்றைவிடத்
தளர்வானது. தன்னிச்சையான எந்த முடிவுறு குறியெழுத்துத் தொகுதியையும் நாம்
அனுமதிக்கிறோம்; மேலும் கட்டுப்பாடுள்ள ஒரு வரையறை $\TMstroke$,
~$\TMblank$ என்ற இரு நாடாக் குறிகளை மட்டும் அனுமதிக்கலாம். பொறி
நாடாவில் ஒரு குறியை எழுதவும் அதே நேரத்தில் நகரவும் நாம்
அனுமதிக்கிறோம்; வேறு வரையறைகள் எழுதுதல் அல்லது நகர்தல் ஆகியவற்றில்
ஒன்றையே அனுமதிக்கின்றன. நாடாவின் தலையை நகர்த்தாமல் எழுதும் வாய்ப்பை
நாம் அனுமதிக்கிறோம்; வேறு வரையறைகள் $\TMstay$ என்ற
``கட்டளையை'' விட்டுவிடுகின்றன. வேறு வகைகளில் நமது வரையறை மேலும்
கட்டுப்பாடுள்ளது. நாடா ஒரு திசையில் மட்டுமே முடிவின்றி நீள்கிறது என்று
நாம் கருதினோம்; வேறு வரையறைகள் நாடா இடப்புறத்திலும் வலப்புறத்திலும்
முடிவின்றி நீள்வதை அனுமதிக்கின்றன. உண்மையில், தனித்தனியான எத்தனை
நாடாக்களையும் அல்லது கட்டங்களின் முடிவிலிக் கட்டத்தளத்தைக்கூட
அனுமதிக்கலாம். டியூரிங் பொறியின் கட்டளைத் தொகுதியை ஒரு நிலைமாற்றுச்
சார்பால் குறிப்பிடுகிறோம்; வேறு வரையறைகள் ஒரு நிலைமாற்று உறவைப்
பயன்படுத்துகின்றன. அங்கு கொடுக்கப்பட்ட எந்தச் சூழலிலும் பொறிக்கு
ஒன்றுக்கு மேற்பட்ட சாத்திய கட்டளைகள் இருக்கும்.

வரையறையின் இந்தக் கடைசித் தளர்வு குறிப்பாக ஆர்வமூட்டுகிறது. நமது
வரையறையில், பொறி நிலை~$q$ இல் குறி~$\sigma$ ஐ வாசிக்கும்போது, புதிய
குறி, நிலை, நாடாவின் தலை இருப்பிடம் ஆகியவை என்னவென்பதை
$\delta(q, \sigma)$ தீர்மானிக்கிறது. ஆனால் நடப்பு நிலை-குறி ஜோடிகள்
$\tuple{q, \sigma}$ க்கும், புதிய நிலை-குறி-திசை மும்மடிகள்
$\tuple{q', \sigma', D}$ க்கும் இடையிலான ஓர் உறவாகக் கட்டளைத்
தொகுதியை அனுமதித்தால், டியூரிங் பொறியின் செயல் தனித்தவாறு
தீர்மானிக்கப்படாமல் இருக்கலாம்---கட்டளை உறவு $\tuple{q, \sigma, q',
\sigma', D}$, $\tuple{q, \sigma, q'', \sigma'', D'}$ ஆகிய இரண்டையும்
கொண்டிருக்கலாம். இந்நேர்வில் நம்மிடம் ஒரு \emph{தீர்மானமற்ற} டியூரிங்
பொறி உள்ளது. கணக்கீட்டுச் சிக்கல்தன்மைக் கோட்பாட்டில் இவை முக்கியப்
பங்கு வகிக்கின்றன.

ஒரு டியூரிங் பொறி எப்போது நிற்கிறது என்பதற்கும் வெவ்வேறு மரபுகள்
உள்ளன: நிலைமாற்றுச் சார்பு வரையறுக்கப்படாதபோது அது நிற்கிறது என்று
நாம் கூறுகிறோம்; வேறு வரையறைகளில் பொறி சிறப்பாக ஒதுக்கப்பட்ட ஒரு
நிற்றல் நிலையில் இருக்க வேண்டும். ஒதுக்கப்பட்ட ஒரு நிற்றல் நிலையைக்
கோருவது டியூரிங் பொறிகளால் எதைக் கணிக்க முடியும் என்பதைப் பாதிக்கும்
ஒரு கட்டுப்பாடு அல்ல என்பதை \olref[hal]{sec} இல் விளக்கியுள்ளோம்.
நமது டியூரிங் பொறிகளின் நாடாக்கள் ஒரு திசையில் மட்டுமே முடிவின்றி
நீள்வதால், ஒரு டியூரிங் பொறியால் ஒரு கட்டளையைச் சரியாக நிறைவேற்ற
முடியாத நேர்வுகள் உள்ளன: அது இடக்கோடிக் கட்டத்தை வாசித்துக்கொண்டு
இடப்புறம் நகர வேண்டியிருந்தால் இவ்வாறு நிகழும். நமது வரையறையின்படி,
``வெளியே விழுவதற்குப்'' பதிலாக அது இருக்கும் இடத்திலேயே தங்குகிறது;
ஆனால் அவ்வாறு நிகழும்போது அது நிற்கும் வகையிலும் நாம் வரையறுத்திருக்கலாம்.
இவ்வரையறையும் சமானமானது: கட்டம்~$0$ இலிருந்து இடப்புறம் நகர முயலும்போது
நிற்கும் டியூரிங் பொறியின் நடத்தையை, ஒவ்வொரு நிலைமாற்றத்தையும்
$\delta(q,\TMendtape) =
\tuple{q',\sigma,\TMleft}$ நீக்குவதன் மூலம் உருவகப்படுத்தலாம்---அப்போது
~$\TMendtape$ இல் இடப்புறம் நகர முயல்வதற்குப் பதிலாகப் பொறி நிற்கும்.
\footnote{இது \emph{முழுமையாகச்} செயல்படாது; ஏனெனில் கட்டம்~$0$
அல்லாத பிற கட்டங்களில் $\TMendtape$ ஐ எழுதுவதையும் வாசிப்பதையும்
எதுவும் தடுக்கவில்லை (\olref[una]{ex:mover} ஐப் பார்க்க). அத்தகைய
நோக்கத்துக்கு அதற்குப் பதிலாகப் பயன்படுத்த இரண்டாவது~$\TMendtape'$
குறியைச் சேர்ப்பதன் மூலம் இதைச் சமாளிக்கலாம்.}

எண்களைக் குறிப்பிடுவதற்கும் (அதனால் ஒரு டியூரிங் பொறி கணிக்கும்
உள்ளீடு-வெளியீட்டுச் சார்பைக் குறிப்பிடுவதற்கும்) வெவ்வேறு வழிகள்
உள்ளன: நாம் ஒருமக் குறிப்பிடலைப் பயன்படுத்துகிறோம்; ஆனால் இருமக்
குறிப்பிடலையும் பயன்படுத்தலாம். இதற்கு $\TMblank$,~$\TMendtape$
ஆகியவற்றுடன் மேலும் இரு குறிகள் தேவை.

இப்போது ஓர் ஆர்வமூட்டும் உண்மை: எந்தச் சார்புகள்
டியூரிங்-கணிக்கத்தக்கவை என்பதில் இம்மாறுபாடுகள் எதுவும் வேறுபாட்டை
ஏற்படுத்துவதில்லை. \emph{ஒரு வரையறையின்படி ஒரு சார்பு
டியூரிங்-கணிக்கத்தக்கது எனில், எல்லா வரையறைகளின்படியும் அது
டியூரிங்-கணிக்கத்தக்கது.}

இதைச் சரிபார்ப்பதன் விவரங்களுக்குள் நாம் செல்ல மாட்டோம். ஒரே ஓர்
எடுத்துக்காட்டு இதோ: இரு திசைகளிலும் முடிவின்றி நீளும் ஒரு நாடாவையோ
பல நாடாக்களையோ அனுமதிப்பதால் கூடுதல் கணிப்புத் திறன் எதையும் நாம்
பெறுவதில்லை. ஏறத்தாழச் சொன்னால், ஒரே ஒரு திசையில் முடிவின்றி நீளும்
நாடாவைக் கொண்ட ஒரு டியூரிங் பொறி, பல நாடாக்களையோ இரு திசைகளிலும்
முடிவின்றி நீளும் நாடாக்களையோ உருவகப்படுத்த முடியும் என்பதே இதற்குக்
காரணம். எடுத்துக்காட்டாக, கூடுதல் நிலைகளையும் கட்டளைகளையும் பயன்படுத்தி,
பல நாடாக்களையோ இரு திசைகளிலும் முடிவின்றி நீளும் நாடாவையோ கொண்ட ஒரு
பொறிக்கான நிரலை, ஒரு திசையில் முடிவின்றி நீளும் ஒரே நாடாவைக் கொண்ட
பொறிக்கான நிரலாக ``மொழிபெயர்க்கலாம்''. மொழிபெயர்க்கப்பட்ட பொறி இரட்டை
எண் கொண்ட கட்டங்களை நாடா~$1$ இன் கட்டங்களுக்காக (அல்லது இரு திசைகளிலும்
முடிவின்றி நீளும் ஒரு நாடாவின் ``நேர்ம'' கட்டங்களுக்காக) பயன்படுத்தலாம்;
ஒற்றை எண் கொண்ட கட்டங்களை நாடா~$2$ இன் கட்டங்களுக்காக (அல்லது
``எதிர்ம'' கட்டங்களுக்காக) பயன்படுத்தலாம்.

\end{document}
